\section{Master certificates and checked selection}
\label{sec:certificates}

\subsection{Why one certificate matters}

A large finitary proof often introduces many local witnesses: a trace, a finite net, branch choices, error allocations, local algebraic identities, a spanning tree, and a list of bounded solutions.  Treating these as unrelated existential quantifiers invites hidden choice and makes the final theorem difficult to audit.  FINER instead assigns every local field a finite radix and packs all fields into one number \(c<M(q)\).

The checker \(\CertCheck(q,c)\) performs four kinds of tests:
\begin{enumerate}[label=(\roman*),leftmargin=2.3em]
\item digit and array range checks;
\item local defining-equation and arithmetic checks;
\item edge checks linking dependent node certificates;
\item terminal checks connecting the decoded observation to the target formula.
\end{enumerate}
The checker is normalized to a Boolean bit and reflected by \cref{thm:boolean-reflection}.  Therefore the semantic-looking statement ``the certificate passes'' is an object-language equation \(\CertCheck(q,c)=1\).

\subsection{Least accepted code}

Let
\[
  f_q(c)=1\dot{-}\CertCheck(q,c).
\]
The checked search module returns
\[
  \Sel(q)=\mu c<M(q)\,[f_q(c)=0].
\]
For a normalized checker bit, \(f_q(c)=0\) iff \(\CertCheck(q,c)=1\).

\begin{theorem}[Selected-certificate theorem]
\label{thm:selected-certificate}
Assume the object theory derives
\begin{equation}
  \exists c<M(q)\;[\CertCheck(q,c)=1].
  \label{eq:accepted-exists}
\end{equation}
Then it derives
\begin{align}
  &\Sel(q)<M(q),\label{eq:selected-bound}\\
  &\CertCheck(q,\Sel(q))=1,\label{eq:selected-accept}\\
  &c<\Sel(q)\;\to\;\CertCheck(q,c)=0.\label{eq:selected-minimal}
\end{align}
\end{theorem}

\begin{proof}
From \eqref{eq:accepted-exists}, existential elimination supplies a witness \(c\) and proofs of \(c<M(q)\) and \(\CertCheck(q,c)=1\).  Bit normalization converts the latter to \(f_q(c)=0\).  Apply \texttt{Search.search\_found}; it gives that the search result is below the bound, is a zero of \(f_q\), and is no larger than the supplied witness.  Convert \(f_q(\Sel(q))=0\) back to \(\CertCheck(q,\Sel(q))=1\).  Minimality is \texttt{Search.search\_minimal} with the same conversion.
\end{proof}

The leastness is not mathematically necessary for soundness, but it has three practical benefits.  It makes outputs deterministic, gives a canonical regression fixture, and permits exact comparison of independently generated certificates.

\subsection{Success and failure evidence}

When a bounded problem is decidable, the framework can expose both outcomes.

\begin{definition}[Dual certificate interface]
A dual interface has a success checker \(C^+\), a failure checker \(C^-\), and object proofs that:
\begin{align*}
 C^+(q,c)=1 &\to \Phi(q),\\
 C^-(q,d)=1 &\to \neg\Phi(q),\\
 \bigl(\exists c<M^+(q)\,C^+(q,c)=1\bigr)
 &\vee
 \bigl(\exists d<M^-(q)\,C^-(q,d)=1\bigr).
\end{align*}
\end{definition}

For a fixed bounded formula, the Boolean reflection construction supplies such an interface.  For a bounded existential, success is a witness and failure is a finite table of rejected candidates, or simply the reflected universal statement that every candidate fails.  For a finite game, success and failure are dual strategy trees.  The important property is not that every theorem has a short negative certificate, but that bounded classical case analysis becomes finite evidence on both branches.

\subsection{Certificate soundness is compositional}

Suppose a graph has nodes \(v_0,\ldots,v_{m-1}\) in a dependency order after demand closure and finite unrolling.  Let \(A_i(q,c)\) be the assertion that node \(i\)'s decoded block is valid.  The master checker is a Boolean conjunction of their reflected bits.  Reflection yields each \(A_i\) from master acceptance.

For every node, provide a local sequent
\begin{equation}
  A_i,\;\{P_j:j\in\operatorname{pred}(i)\},\;\Gamma_i
  \;\vdash[\PRAAx]\;P_i,
  \label{eq:local-node-sequent}
\end{equation}
where \(P_j\) are predecessor postconditions and \(\Gamma_i\) are residual assumptions used locally.  Repeated simultaneous cut along the finite order yields the terminal postcondition.  At no point does the compiler need to inspect the semantic meaning of \(P_i\); it only substitutes terms and composes derivations.

This is the proof-theoretic reason for using one master certificate.  Its accepted equation is a single reusable major premise from which all node facts are projected.

\subsection{Bounds and enormous finite spaces}

The certificate space may be astronomically large.  That does not invalidate the object theorem.  Primitive recursive arithmetic concerns total finite computations, not feasible ones.  Nevertheless, practical formalization benefits from factorization:

\begin{itemize}[leftmargin=2em]
\item deterministic fields should be recomputed rather than stored;
\item independent blocks should use local searches whose results are then packed;
\item bounds should be stated symbolically rather than expanded as numerals;
\item proof checking should use reflected local equations, not execute the entire search in Lean's kernel.
\end{itemize}

The role of least search is logical witness extraction.  An external producer may directly supply a compact certificate for testing and replay; the object theorem remains valid because it proves the least accepted certificate exists whenever any accepted certificate exists.
